% GENERATED FILE. Edit canonical lessons, then run npm run generate:books.
% canonical-source-hash: fnv1a64:f98aae4e30e500ea
% canonical-lessons: JA-C506-kyougi, JA-C506-senshuken, JA-C506-shouri, JA-C506-haiboku, JA-C506-makasu

\chapter{Winning and Losing}
\label{ch:ja-winning-and-losing}
\hlchaptermodality{506}

\begin{chapteropening}
\textbf{By the end of this chapter:} \emph{I can talk about contests, wins and defeats.}

Complete the last lesson of chapter 506: I can talk about contests, wins and defeats.
\end{chapteropening}

\section[\texorpdfstring{\ja{きょうぎ} (kyōgi)}{きょうぎ (kyōgi)}]{\ja{きょうぎ} (kyōgi) — a sport; a competition}

\label{lesson:JA-C506-kyougi}



\begin{quote}
Before the new one: say the Japanese for a bow (for arrows), then the Japanese for to dodge.
\end{quote}

\subsection*{You'll want to know: \ja{きょうぎ}}
\textbf{\ja{きょうぎ}} — \emph{kyōgi} — \textquotedblleft{}a sport; a competition\textquotedblright{}.

\begin{quote}
\textbf{\ja{どの} \ja{きょうぎが} \ja{すきですか。}} — \emph{dono kyōgi ga suki desu ka} — \textquotedblleft{}Which sport do you like?\textquotedblright{}
\end{quote}

\begin{cousinweb}
\textbf{\ja{きょうぎ}} (\emph{kyōgi}) is Sino-Japanese (kango): two Chinese characters read with their on'yomi, readings borrowed from Chinese. They are \emph{kyō}, \textquotedblleft{}emulate\textquotedblright{} (Middle Chinese \emph{gjaengH}) and \emph{gi}, \textquotedblleft{}skill\textquotedblright{} (\emph{gjeX}) (KANJIDIC2; Middle Chinese in Baxter's transcription).
\end{cousinweb}

\begin{grammarlens}[title={What you've built}]
One more word for describing what you lived and felt.
\end{grammarlens}

\subsection*{Your turn}
\begin{itemize}
\raggedright
  \item \emph{Say it:} \emph{kyōgi}
  \item \emph{Say it:} \emph{kyōgi}, once more
  \item \emph{Say it:} \emph{dono kyōgi ga suki desu ka} — Which sport do you like?
  \item \emph{Recall:} say the Japanese for a bow (for arrows), then the Japanese for to dodge, then say \emph{kyōgi} again
\end{itemize}

\begin{tcolorbox}[breakable,colback=teal!4,colframe=teal!35!black,title={Before you move on}]
What does \ja{きょうぎ} mean? (\textquotedblleft{}A sport; a competition\textquotedblright{}.) Say it once more.
\end{tcolorbox}
\section[\texorpdfstring{\ja{せんしゅけん} (senshuken)}{せんしゅけん (senshuken)}]{\ja{せんしゅけん} (senshuken) — a championship}

\label{lesson:JA-C506-senshuken}



\begin{quote}
Before the new one: say the Japanese for to dodge, then the Japanese for a sport.
\end{quote}

\subsection*{You'll want to know: \ja{せんしゅけん}}
\textbf{\ja{せんしゅけん}} — \emph{senshuken} — \textquotedblleft{}a championship\textquotedblright{}.

\begin{quote}
\textbf{\ja{あには} \ja{すいえいの} \ja{せんしゅけんで} \ja{およぎます。}} — \emph{ani wa suiei no senshuken de oyogimasu} — \textquotedblleft{}My elder brother swims in the swimming championship.\textquotedblright{}
\end{quote}

\begin{cousinweb}
\textbf{\ja{せんしゅけん}} (\emph{senshuken}) is Sino-Japanese (kango): three Chinese characters read with their on'yomi, readings borrowed from Chinese. They are \emph{sen}, \textquotedblleft{}elect\textquotedblright{} (Middle Chinese \emph{sjwenX}), \emph{shu}, \textquotedblleft{}hand\textquotedblright{} (\emph{syuwX}) and \emph{ken}, \textquotedblleft{}authority\textquotedblright{} (\emph{gjwen}) (KANJIDIC2; Middle Chinese in Baxter's transcription).
\end{cousinweb}

\begin{grammarlens}[title={What you've built}]
One more word for describing what you lived and felt.
\end{grammarlens}

\subsection*{Your turn}
\begin{itemize}
\raggedright
  \item \emph{Say it:} \emph{senshuken}
  \item \emph{Say it:} \emph{senshuken}, once more
  \item \emph{Say it:} \emph{ani wa suiei no senshuken de oyogimasu} — My elder brother swims in the swimming championship.
  \item \emph{Recall:} say the Japanese for to dodge, then the Japanese for a sport, then say \emph{senshuken} again
\end{itemize}

\begin{tcolorbox}[breakable,colback=teal!4,colframe=teal!35!black,title={Before you move on}]
What does \ja{せんしゅけん} mean? (\textquotedblleft{}A championship\textquotedblright{}.) Say it once more.
\end{tcolorbox}
\section[\texorpdfstring{\ja{しょうり} (shōri)}{しょうり (shōri)}]{\ja{しょうり} (shōri) — a victory}

\label{lesson:JA-C506-shouri}



\begin{quote}
Before the new one: say the Japanese for a sport, then the Japanese for a championship.
\end{quote}

\subsection*{You'll want to know: \ja{しょうり}}
\textbf{\ja{しょうり}} — \emph{shōri} — \textquotedblleft{}a victory\textquotedblright{}.

\begin{quote}
\textbf{\ja{みんなで} \ja{しょうりを} \ja{よろこびました。}} — \emph{minna de shōri o yorokobimashita} — \textquotedblleft{}We all celebrated the victory.\textquotedblright{}
\end{quote}

\begin{cousinweb}
\textbf{\ja{しょうり}} (\emph{shōri}) is Sino-Japanese (kango): two Chinese characters read with their on'yomi, readings borrowed from Chinese. They are \emph{shō}, \textquotedblleft{}victory\textquotedblright{} (Middle Chinese \emph{syingH}) and \emph{ri}, \textquotedblleft{}profit\textquotedblright{} (\emph{lijH}) (KANJIDIC2; Middle Chinese in Baxter's transcription).
\end{cousinweb}

\begin{grammarlens}[title={What you've built}]
One more word for describing what you lived and felt.
\end{grammarlens}

\subsection*{Your turn}
\begin{itemize}
\raggedright
  \item \emph{Say it:} \emph{shōri}
  \item \emph{Say it:} \emph{shōri}, once more
  \item \emph{Say it:} \emph{minna de shōri o yorokobimashita} — We all celebrated the victory.
  \item \emph{Recall:} say the Japanese for a sport, then the Japanese for a championship, then say \emph{shōri} again
\end{itemize}

\begin{tcolorbox}[breakable,colback=teal!4,colframe=teal!35!black,title={Before you move on}]
What does \ja{しょうり} mean? (\textquotedblleft{}A victory\textquotedblright{}.) Say it once more.
\end{tcolorbox}
\section[\texorpdfstring{\ja{はいぼく} (haiboku)}{はいぼく (haiboku)}]{\ja{はいぼく} (haiboku) — a defeat}

\label{lesson:JA-C506-haiboku}



\begin{quote}
Before the new one: say the Japanese for a championship, then the Japanese for a victory.
\end{quote}

\subsection*{You'll want to know: \ja{はいぼく}}
\textbf{\ja{はいぼく}} — \emph{haiboku} — \textquotedblleft{}a defeat\textquotedblright{}.

\begin{quote}
\textbf{\ja{はいぼくは} \ja{くやしいです。}} — \emph{haiboku wa kuyashii desu} — \textquotedblleft{}Defeat is bitter.\textquotedblright{}
\end{quote}

\begin{cousinweb}
\textbf{\ja{はいぼく}} (\emph{haiboku}) is Sino-Japanese (kango): two Chinese characters read with their on'yomi, readings borrowed from Chinese. They are \emph{hai}, \textquotedblleft{}failure\textquotedblright{} (Middle Chinese \emph{baejH}) and \emph{boku}, \textquotedblleft{}north\textquotedblright{} (\emph{pok}) (KANJIDIC2; Middle Chinese in Baxter's transcription).
\end{cousinweb}

\begin{grammarlens}[title={What you've built}]
One more word for describing what you lived and felt.
\end{grammarlens}

\subsection*{Your turn}
\begin{itemize}
\raggedright
  \item \emph{Say it:} \emph{haiboku}
  \item \emph{Say it:} \emph{haiboku}, once more
  \item \emph{Say it:} \emph{haiboku wa kuyashii desu} — Defeat is bitter.
  \item \emph{Recall:} say the Japanese for a championship, then the Japanese for a victory, then say \emph{haiboku} again
\end{itemize}

\begin{tcolorbox}[breakable,colback=teal!4,colframe=teal!35!black,title={Before you move on}]
What does \ja{はいぼく} mean? (\textquotedblleft{}A defeat\textquotedblright{}.) Say it once more.
\end{tcolorbox}
\section[\texorpdfstring{\ja{まかす} (makasu)}{まかす (makasu)}]{\ja{まかす} (makasu) — to defeat; to beat}

\label{lesson:JA-C506-makasu}



\begin{quote}
Before the new one: say the Japanese for a victory, then the Japanese for a defeat.
\end{quote}

\subsection*{You'll want to know: \ja{まかす}}
\textbf{\ja{まかす}} — \emph{makasu} — \textquotedblleft{}to defeat; to beat\textquotedblright{}.

\begin{quote}
\textbf{\ja{おとうとが} \ja{しょうぎで} \ja{ちちを} \ja{まかしました。}} — \emph{otōto ga shōgi de chichi o makashimashita} — \textquotedblleft{}My younger brother beat my father at shogi.\textquotedblright{}
\end{quote}

Polite \textbf{\ja{まかします}}, \emph{makashimasu}; a godan verb. Past (polite): \textbf{\ja{いもうとが} \ja{あにを} \ja{まかしました。}}, \emph{imōto ga ani o makashimashita}.

\begin{cousinweb}
\textbf{\ja{まかす}} (\emph{makasu}) is a native Japanese word (wago), inherited rather than borrowed from Chinese; its kanji was attached later as a kun'yomi, and its older origin is not settled (JMdict, KANJIDIC2).
\end{cousinweb}

\begin{grammarlens}[title={What you've built}]
One more word for describing what you lived and felt.
\end{grammarlens}

\subsection*{Your turn}
\begin{itemize}
\raggedright
  \item \emph{Say it:} \emph{makasu}
  \item \emph{Say it:} \emph{makasu}, once more
  \item \emph{Say it:} \emph{otōto ga shōgi de chichi o makashimashita} — My younger brother beat my father at shogi.
  \item \emph{Recall:} say the Japanese for a victory, then the Japanese for a defeat, then say \emph{makasu} again
\end{itemize}

\begin{tcolorbox}[breakable,colback=teal!4,colframe=teal!35!black,title={Before you move on}]
What does \ja{まかす} mean? (\textquotedblleft{}To defeat; to beat\textquotedblright{}.) Say it once more.
\end{tcolorbox}
